\newglossaryentry{anamnese}{
    name={Anamnese},
    description={levantamento do histórico do paciente, de suas condições de saúde e dos sinais relatados pelo tutor, utilizado para orientar a avaliação clínica}
}

\newglossaryentry{armazenamento-chave-valor}{
    name={Armazenamento chave-valor},
    description={forma de armazenamento na qual cada valor é associado a uma chave que permite sua identificação e recuperação}
}

\newglossaryentry{atomicidade}{
    name={Atomicidade},
    description={propriedade de uma transação que assegura a efetivação de todas as suas operações ou a reversão do conjunto em caso de falha}
}

\newglossaryentry{autenticacao}{
    name={Autenticação},
    description={processo de verificação da identidade de um usuário ou sistema por meio de credenciais}
}

\newglossaryentry{back-end}{
    name={Back-end},
    description={camada da aplicação responsável pelo processamento de requisições, pela aplicação de regras de negócio e pelo acesso aos dados mantidos no servidor}
}

\newglossaryentry{banco-relacional}{
    name={Banco de dados relacional},
    description={banco de dados que organiza informações em tabelas compostas por linhas e colunas, com relacionamentos representados por chaves}
}

\newglossaryentry{chave-estrangeira}{
    name={Chave estrangeira},
    description={campo ou conjunto de campos que referencia uma chave de outra tabela, ou da própria tabela, estabelecendo um relacionamento entre registros}
}

\newglossaryentry{chave-primaria}{
    name={Chave primária},
    description={campo ou conjunto de campos que identifica de forma única cada registro de uma tabela}
}

\newglossaryentry{cliente-servidor}{
    name={Cliente-servidor},
    description={modelo de arquitetura no qual o cliente solicita serviços ou recursos e o servidor recebe e processa essas solicitações, retornando respostas}
}

\newglossaryentry{conflito-sincronizacao}{
    name={Conflito de sincronização},
    description={situação em que alterações concorrentes de um mesmo registro precisam ser conciliadas para determinar o estado dos dados}
}

\newglossaryentry{endpoint}{
    name={Endpoint},
    description={ponto de acesso de uma interface de programação de aplicações, identificado por um endereço ao qual são dirigidas requisições}
}

\newglossaryentry{exclusao-logica}{
    name={Exclusão lógica},
    description={procedimento que marca um registro como excluído sem removê-lo fisicamente de imediato do armazenamento}
}

\newglossaryentry{fila-sincronizacao}{
    name={Fila de sincronização},
    description={estrutura que mantém operações pendentes para posterior transmissão e processamento durante a sincronização de dados}
}

\newglossaryentry{front-end}{
    name={Front-end},
    description={camada da aplicação que apresenta a interface ao usuário e recebe suas interações}
}

\newglossaryentry{hash}{
    name={Hash},
    description={resultado de uma função que transforma dados de entrada em um valor de tamanho definido; no armazenamento de senhas, devem ser empregados mecanismos específicos de derivação com proteção contra tentativas de descoberta}
}

\newglossaryentry{interoperabilidade}{
    name={Interoperabilidade},
    description={capacidade de sistemas ou componentes trocarem informações e utilizarem os dados recebidos conforme formatos e contratos compatíveis}
}

\newglossaryentry{last-write-wins}{
    name={Last-Write-Wins},
    description={estratégia de resolução de conflitos que seleciona a versão considerada mais recente de um dado com base na ordenação das escritas, geralmente determinada por marcas temporais}
}

\newglossaryentry{metadados}{
    name={Metadados},
    description={dados que descrevem ou contextualizam outros dados, como identificadores, datas de alteração e informações de controle da sincronização}
}

\newglossaryentry{offline-first}{
    name={Offline-first},
    description={abordagem de desenvolvimento que prioriza o funcionamento da aplicação com dados locais, permitindo executar suas funções previstas sem conexão e sincronizar alterações quando houver conectividade}
}

\newglossaryentry{persistencia-local}{
    name={Persistência local},
    description={armazenamento de dados no dispositivo do usuário de forma que permaneçam disponíveis entre execuções da aplicação, inclusive sem acesso à rede}
}

\newglossaryentry{prontuario-clinico-veterinario}{
    name={Prontuário clínico veterinário},
    description={conjunto organizado de registros referentes à identificação do paciente animal, ao histórico de saúde e aos atendimentos, incluindo avaliações, condutas e documentos clínicos}
}

\newglossaryentry{serializacao}{
    name={Serialização},
    description={processo de conversão de dados estruturados em uma representação adequada ao armazenamento ou à transmissão, como texto no formato JSON}
}

\newglossaryentry{sincronizacao-assincrona}{
    name={Sincronização assíncrona},
    description={processo de troca e conciliação de dados executado sem exigir que a interação principal do usuário aguarde sua conclusão}
}

\newglossaryentry{sincronizacao-incremental}{
    name={Sincronização incremental},
    description={processo que transmite as alterações identificadas desde um marco de sincronização anterior, em vez de reenviar integralmente todos os registros}
}

\newglossaryentry{timestamp}{
    name={Timestamp},
    description={marca temporal que representa um instante e pode ser utilizada para registrar alterações e comparar versões de dados}
}

\newglossaryentry{token-autenticacao}{
    name={Token de autenticação},
    description={credencial utilizada para comprovar uma autenticação prévia em requisições ao sistema, sujeita a regras de validação e validade}
}

\newglossaryentry{transacao}{
    name={Transação},
    description={conjunto de operações de banco de dados tratado como uma unidade de processamento, cuja confirmação ou reversão preserva as propriedades estabelecidas pelo sistema}
}

\newglossaryentry{tutor}{
    name={Tutor},
    description={pessoa responsável pelo animal e por fornecer informações e acompanhar os cuidados relacionados ao atendimento veterinário}
}

